\section{Rigidity and biinterpretability}
\label{sec:rigidity}

We now apply the recovery construction to the full degree
structure.

\begin{proof}[Proof of Theorem~\ref{thm:rigidity}]
Fix an arbitrary order automorphism $\pi$ and take the
Borel map $F$ given by Lemma~\ref{lem:lift}. Fix an
irrational $t\in(0,1)$. Proposition~\ref{prop:recovery}
gives a comeager set $G_t$ of pairs $(u,v)$ for which
\eqref{eq:four-values} holds with a suitable rational
$\delta\in(0,1)$.

For such a pair, put
$\mathbf b=d(t)\vee d(u)\vee d(v)$.
Each argument in the four values on the right of
\eqref{eq:four-values} is a rational linear combination
of $t,u,v,1$. Lemma~\ref{lem:cuts} bounds its degree by
$\mathbf b$, and Equation~\eqref{eq:representation}
bounds the degree of its $F$-value by $\pi(\mathbf b)$.
The same upper bound holds for the finite join of the
four values. Thus
\[
d(t)\leT \pi\bigl(d(t)\vee d(u)\vee d(v)\bigr)
\qquad((u,v)\in G_t).
\]
Applying the inverse order automorphism yields
\begin{equation}\label{eq:generic-bound}
\pi^{-1}(d(t))\leT d(t)\vee d(u)\vee d(v)
\qquad((u,v)\in G_t).
\end{equation}

Choose one fixed set $Y$ representing $\pi^{-1}(d(t))$.
If $Y\not\leT C(t)$, Lemma~\ref{lem:avoidance} says that
\[
\{(u,v):Y\leT C(t)\oplus C(u)\oplus C(v)\}
\]
is meager. But \eqref{eq:generic-bound} puts the
comeager set $G_t$ inside this set, contradicting the
Baire Category Theorem for $\R^2$. Consequently
\[
\pi^{-1}(d(t))\leT d(t).
\]
The set $Y$ is fixed before the pair varies. No common
index for the reductions in \eqref{eq:generic-bound}
is required.

The irrational $t$ was arbitrary. By Lemma~\ref{lem:cuts},
its degrees cover every noncomputable degree, and the
least degree is fixed. We have therefore proved
\begin{equation}\label{eq:one-sided}
\pi^{-1}(\mathbf a)\leT\mathbf a
\qquad(\mathbf a\in\DT)
\end{equation}
for every order automorphism $\pi$.
Apply this result to the automorphism $\pi^{-1}$ to
obtain $\pi(\mathbf a)\leT\mathbf a$. Applying $\pi$
to \eqref{eq:one-sided} gives
$\mathbf a\leT\pi(\mathbf a)$. Antisymmetry now gives
$\pi(\mathbf a)=\mathbf a$ for every degree.
\end{proof}

The second application uses Theorem~\ref{thm:representation}
for the inverse degree automorphism. It does not require
an inverse of $F$, which may have nontrivial fibers.

\begin{corollary}\label{cor:biinterpretability}
The Turing-degree ordering is biinterpretable without parameters
with full second-order arithmetic, in the sense of
\cite[Definition~7.5.1]{SlamanWoodin2005}.
\end{corollary}

\begin{proof}
Slaman and Woodin prove that this property is equivalent to
rigidity \cite[Theorem~7.5.3]{SlamanWoodin2005}.
Apply Theorem~\ref{thm:rigidity}.
\end{proof}
